% ECONOMICS_PROBLEM: {"id": "C72-620", "title": "Structural classes with polynomial-time exact bimatrix Nash computation", "jel_code": "C72", "source_id": "OP-0135"}
% !TeX program = xelatex
\documentclass[11pt,a4paper]{article}
\usepackage[margin=23mm]{geometry}
\usepackage{fontspec}
\setmainfont{Latin Modern Roman}
\usepackage{amsmath,amssymb,mathtools}
\usepackage{xcolor,enumitem,booktabs,longtable,array,xurl,hyperref}
\definecolor{muted}{HTML}{53636C}
\hypersetup{colorlinks=true,urlcolor=blue,linkcolor=blue}
\setlength{\parindent}{0pt}
\setlength{\parskip}{5pt}
\newcommand{\R}{\mathbb R}
\newcommand{\E}{\mathbb E}
\newcommand{\N}{\mathbb N}
\newcommand{\ind}{\mathbf 1}
\newcommand{\Var}{\operatorname{Var}}
\newcommand{\Cov}{\operatorname{Cov}}
\newcommand{\supp}{\operatorname{supp}}
\DeclareMathOperator*{\argmin}{arg\,min}
\DeclareMathOperator*{\argmax}{arg\,max}
\newcommand{\entrymeta}[3]{{\sffamily\small\color{muted}\textbf{\texttt{#1}}\quad|\quad #2\par #3\par}}
\newcommand{\Problemref}[1]{\texttt{#1}}
\newcommand{\JELref}[2]{\texttt{#2} (JEL \texttt{#1})}
\begin{document}
\section*{Structural classes with polynomial-time exact bimatrix Nash computation}
\textbf{Problem ID:} \texttt{C72-620}\quad \textbf{JEL:} \texttt{C72}\par
% BEGIN ECONOMICS STATEMENT
\label{problem:OP-0135}
% GAME_THEORY_PROBLEM: {"local_id": "GT-005", "original_number": 5, "kind": "R", "entry_kind": "statement_only", "published": false, "review_required": true, "category": "game_theory"}
\label{game:GT-005}
{\sffamily\small\color{muted}
\textbf{GT-005} | Algorithmic equilibrium and complexity\\
\textbf{R} | Classification and algorithm-design agenda.
\par}
\subsection*{Definitions and mathematical statement}
Let $\mathcal B$ be the set of explicitly encoded rational bimatrix games. For $G=(A,B)$, let $\operatorname{NE}(G)$ be the set of $(x,y)\in\Delta_m\times\Delta_n$ satisfying
\[
 (Ay)_i\leq x^{\mathsf T}Ay\quad(1\leq i\leq m),\qquad
 (x^{\mathsf T}B)_j\leq x^{\mathsf T}By\quad(1\leq j\leq n).
\]
A structural family is specified by a mathematical predicate $P$ on the game's payoffs and, if necessary, by an explicitly encoded witness $w$ for the structural promise. The target is to identify a substantively motivated family $\mathcal C=\{(G,w):P(G,w)\}$ and an algorithm $\mathcal A$ for which
\[
 \forall(G,w)\in\mathcal C,\quad
 \mathcal A(G,w)\in\operatorname{NE}(G)\cap(\mathbb Q^m\times\mathbb Q^n),\qquad
 T_{\mathcal A}(G,w)\leq p(|G|+|w|)
\]
for some polynomial $p$. The family, its encoding, recognition or promise convention, and its relation to already tractable families must be part of the result.
\subsection*{Short English statement}
Find structural assumptions that make exact two-player Nash equilibrium computation efficient.
\subsection*{Variants and scope boundaries}
There is no well-defined ``largest tractable class'' without a chosen ordering and restrictions on family descriptions. Classes obtained by simply enumerating preselected easy instances do not answer the substantive classification agenda. Approximate equilibria, fixed action counts, and exponentially large normal forms are distinct restrictions.
\subsection*{Source relationship and status}
This preserves the source's request for new tractable classes while removing the unsupported implication that it poses one definite yes/no conjecture. Exact bimatrix equilibria admit rational representations; this does not imply polynomial-time computation.
\subsection*{References}
\begingroup\footnotesize\raggedright
{\footnotesize\raggedright
\textbf{[S1]} Hanyu Li, Wenhan Huang, Zhijian Duan, David Henry Mguni, Kun Shao, Jun Wang, and Xiaotie Deng (2023). \emph{A survey on algorithms for Nash equilibria in finite normal-form games}. Locator: concluding ``Open problems and future directions,'' theoretical item 4. \url{https://arxiv.org/pdf/2312.11063}\par}
\par\endgroup

\subsection*{Common conventions: Source mathematical conventions}

\textbf{Finite games.} For a finite set $A$, $\Delta(A)$ is its probability
simplex. Players choose independent mixed strategies in Nash equilibrium;
correlation is allowed only when explicitly part of the solution concept.
In a game with payoff functions $u_i$ and strategy profile $x$, an additive
$\varepsilon$ Nash equilibrium satisfies
\[
 u_i(x)\geq u_i(x_i',x_{-i})-\varepsilon
 \quad\text{for every player $i$ and every admissible deviation $x_i'$}.
\]
For numerical approximation, payoffs are normalized as locally specified.
An $\varepsilon$ payoff residual is distinct from distance at most
$\varepsilon$ to the set of exact equilibria. Point convergence, convergence
of empirical play, and convergence in distance to an equilibrium set are
also distinct.

\textbf{Computational representations.} Unless stated otherwise, a complexity
question takes rational data with binary encoding; polynomial time is measured
in total bit length. A fixed-accuracy polynomial algorithm is not necessarily
polynomial in $1/\varepsilon$, and neither is necessarily polynomial in
$\log(1/\varepsilon)$. Succinct games and value, demand, or sampling oracles
must declare their interfaces. Exact real solutions require an explicit
representation; an unspecified list of arbitrary real numbers is not a
finite computational output. A claimed algorithmic barrier is meaningful
only for its specified model and reductions.

\textbf{Dynamic games.} Histories, information, observation, and allowable
randomization are part of the model. A stationary strategy depends only on
the current state; a positional strategy in a deterministic graph game
chooses one action at each controlled vertex. Nash equilibrium at an initial
state and equilibrium after every history are different requirements.
An expectation of a pathwise $\liminf$ payoff, a $\liminf$ of expected averages,
a limiting expected average, and a uniform finite-horizon equilibrium are
different criteria. None is silently substituted for another.

\textbf{Allocation and mechanisms.} A complete allocation assigns every item
exactly once unless otherwise stated. Zero-valued goods, negative values,
capacity constraints, feasibility, lotteries, and copies of goods affect
fairness definitions and are specified locally. Dominant-strategy incentive
compatibility, Bayesian incentive compatibility, universal truthfulness,
and truthfulness in expectation impose different inequalities. Whenever
an objective ratio has zero denominator, its convention is stated or those
instances are excluded.

\textbf{Infinite-dimensional models.} Expectations, integrals, stochastic
equations, and optimization objectives require their local regularity,
integrability, filtrations, and admissible domains. A backward--forward
mean-field system is a boundary-value problem; it is not an initial-value
semiflow without an additional construction. Long-time, large-population,
and vanishing-noise limits require an order of limits and a convergence
topology. If a minimum need not be attained, the objective is its infimum
or supremum and the approximation requirement is stated separately.
% END ECONOMICS STATEMENT
\end{document}
